\section{繰り込み}

\noindent\textbf{この章の中心命題}
\[
\boxed{
\mathcal{L}(m_0,e_0,\psi_0,A_0)
=
\mathcal{L}_{\mathrm{ren}}(m,e,\psi,A)
+\mathcal{L}_{\mathrm{ct}},
\qquad
\mu\frac{de}{d\mu}=\beta(e)
}
\]
相互作用する場の量子論では、短距離の仮想過程から発散が現れる。繰り込みとは、測定される質量・電荷を基準に理論を書き直し、有限な予言を取り出す手続きである。

QED のラグランジアンを得たことで、電子と光子の放出・吸収を計算できるようになった。しかし相互作用を摂動展開すると、単に「光子を一つ交換する」過程だけでなく、途中で仮想的な電子・陽電子・光子が現れて再び消える過程も含まれる。これらは実験で直接観測される粒子ではないが、散乱振幅やエネルギー準位には寄与する。

\begin{experimentalfact}{ラムシフトと異常磁気モーメント}
\factitem{操作}{水素原子の高精度分光により $2S_{1/2}$ と $2P_{1/2}$ の準位差を測定し、また磁場中の電子の歳差運動から磁気モーメントを測定する。}
\factitem{前提}{ディラック方程式だけを用いると、水素原子の $2S_{1/2}$ と $2P_{1/2}$ は同じエネルギーをもち、点状電子の磁気モーメントは $g=2$ に対応する。}
\factitem{結果}{$2S_{1/2}$ と $2P_{1/2}$ には小さいが測定可能な差があり、電子の
\[
a_e:=\frac{g_e-2}{2}
\]
も $0$ ではない。}
\factitem{示唆}{電磁場の量子ゆらぎや仮想粒子の寄与は、測定可能な物理量に現れる。}
\end{experimentalfact}

この実験事実は、QED の相互作用を単なる古典的な電磁場との結合としてではなく、量子場の相互作用として計算する必要があることを示している。

\subsection{摂動展開と発散の出現}
相互作用項
\[
\mathcal{L}_{\mathrm{int}}
=
-e\bar{\psi}\gamma^\mu A_\mu\psi
\]
を小さい補正として扱うと、散乱振幅やエネルギー準位は $e$ のべき級数として計算される。最低次では、電子が光子を一つ放出・吸収する過程や、二つの電子が光子を一つ交換する過程が現れる。

次の次数では、閉じたループをもつ図が現れる。たとえば電子が光子を放出し、それを再び吸収する自己エネルギー補正は、形式的には
\[
\Sigma(p)
\sim
e^2
\int \frac{d^4k}{(2\pi)^4}
\gamma^\mu
\frac{\gamma^\rho(p-k)_\rho+m}{(p-k)^2-m^2}
\gamma_\mu
\frac{1}{k^2}
\]
のような積分を含む。ここで $p$ は外から見た電子の運動量、$k$ はループ内を流れる仮想光子の運動量である。

問題は、$k$ が非常に大きい領域、すなわち極めて短い距離の過程まで積分すると、この積分が発散することである。これは、理論が短距離で無限大の補正を予言しているというより、質量や電荷を「相互作用を含まない裸の量」として扱ったままでは、実験で測られる量と対応しないことを示している。

\subsection{裸の量と測定される量}
ラグランジアンに最初に現れる質量や電荷を、裸の質量 $m_0$、裸の電荷 $e_0$ と呼ぶ。また場そのものも裸の場 $\psi_0,A_{0\mu}$ と書く。繰り込みでは、これらを測定される量と補正項に分ける。
\[
\psi_0=Z_2^{1/2}\psi,
\qquad
A_{0\mu}=Z_3^{1/2}A_\mu,
\qquad
m_0=m+\delta m,
\qquad
e_0=Z_e e.
\]
ここで $m,e,\psi,A_\mu$ は繰り込まれた量であり、$\delta m,Z_2,Z_3,Z_e$ は発散を吸収するための係数である。

この書き換えにより、ラグランジアンは
\[
\mathcal{L}(m_0,e_0,\psi_0,A_0)
=
\mathcal{L}_{\mathrm{ren}}(m,e,\psi,A)
+\mathcal{L}_{\mathrm{ct}}
\]
と分かれる。$\mathcal{L}_{\mathrm{ct}}$ は反項と呼ばれ、ループ積分から出る発散を打ち消す役割をもつ。

重要なのは、反項を自由に選ぶのではなく、測定される物理量を基準に固定することである。たとえば電子の質量 $m$ は、電子の伝播関数の極が
\[
p^2=m^2
\]
に来るように定義する。電荷 $e$ は、十分低エネルギーで測ったクーロン力、すなわち長距離での電磁相互作用の強さによって定義する。このように、理論中の記号を測定手順と結び付けることで、発散を含まない有限な予言が得られる。

\subsection{繰り込み可能性}
繰り込みがうまくいくためには、発散を吸収する反項が、もとのラグランジアンと同じ形の項だけで足りる必要がある。QED の場合、必要な反項は
\[
\bar{\psi}i\gamma^\mu\partial_\mu\psi,
\qquad
m\bar{\psi}\psi,
\qquad
F_{\mu\nu}F^{\mu\nu},
\qquad
e\bar{\psi}\gamma^\mu A_\mu\psi
\]
の係数を修正する形で書ける。つまり、新しい種類の相互作用を無限に追加しなくても、有限個の測定値を入力すれば、他の物理量を予言できる。

この性質を繰り込み可能性という。繰り込み可能性は単なる計算上の都合ではなく、短距離の詳細を知らなくても、低エネルギーでの物理を有限個のパラメータで記述できることを意味している。

\subsection{走る結合定数}
繰り込みでは、どのエネルギー尺度で電荷を定義するかを指定する必要がある。その尺度を $\mu$ と書くと、繰り込まれた電荷は一般に $e=e(\mu)$ として $\mu$ に依存する。この依存性は
\[
\mu\frac{de}{d\mu}=\beta(e)
\]
で表される。QED で一種類の荷電ディラック粒子だけを考えると、最低次では
\[
\beta(e)=\frac{e^3}{12\pi^2}+\cdots
\]
である。したがって、QED の有効な電荷は高いエネルギー、すなわち短い距離でわずかに大きくなる。

ここで重要なのは、「電荷が変わる」という言い方が、同じ実験を何度も行うと電荷が時間的に変動するという意味ではないことである。測定で探る距離スケールが変わると、真空中の仮想電子・陽電子対による遮蔽のされ方が変わり、その結果として有効な相互作用の強さが変わって見える、という意味である。

\subsection{繰り込みの物理的意味}
繰り込みは、無限大を形式的に隠すための技巧ではない。物理的には、測定装置が有限の距離スケールでしか系を分解できないことを理論に反映させる手続きである。非常に短い距離の自由度は、低エネルギーの観測では直接見えず、質量・電荷・場の規格化などの値にまとめて吸収される。

したがって、繰り込みを通じて得られる見方は次のようにまとめられる。相互作用する場の量子論では、ラグランジアンに書いたパラメータそのものではなく、どのスケールでどの物理量を測って定義したかが基本になる。そして、その入力から別のスケールや別の過程での観測量を計算することが、理論の予言である。
